% docs/slides/preamble.tex -- 五份 deck 共享
\documentclass[aspectratio=169,11pt]{beamer}

% ==== 中文与字体 ====
\usepackage[UTF8,fontset=none]{ctex}
\usepackage{fontspec}

\IfFontExistsTF{PingFang SC}{%
  \setCJKmainfont{PingFang SC}[BoldFont=PingFang SC,ItalicFont=PingFang SC,AutoFakeSlant=0.15]
  \setCJKsansfont{PingFang SC}[ItalicFont=PingFang SC,AutoFakeSlant=0.15]
  \setCJKmonofont{PingFang SC}
}{%
  \IfFontExistsTF{Source Han Sans SC}{%
    \setCJKmainfont{Source Han Sans SC}
    \setCJKsansfont{Source Han Sans SC}
  }{%
    \IfFontExistsTF{Noto Sans CJK SC}{%
      \setCJKmainfont{Noto Sans CJK SC}
      \setCJKsansfont{Noto Sans CJK SC}
    }{%
      \setCJKmainfont{Heiti SC}
      \setCJKsansfont{Heiti SC}
    }
  }
}

\IfFontExistsTF{Helvetica Neue}{\setmainfont{Helvetica Neue}}{%
  \IfFontExistsTF{Inter}{\setmainfont{Inter}}{}}
\IfFontExistsTF{JetBrains Mono}{\setmonofont{JetBrains Mono}[Scale=0.85]}%
  {\setmonofont{Menlo}[Scale=0.85]}

% ==== 颜色 ====
\usepackage{xcolor}
\definecolor{cMain}{HTML}{1F4E79}
\definecolor{cAccent}{HTML}{E07A1F}
\definecolor{cGray}{HTML}{595959}
\definecolor{cMute}{HTML}{F4F4F4}
\definecolor{cOK}{HTML}{2E7D32}
\definecolor{cBad}{HTML}{B23A3A}
\definecolor{cWarn}{HTML}{C8A300}

% ==== 主题 ====
\usepackage{tcolorbox}
\usetheme{metropolis}
\metroset{progressbar=frametitle,numbering=fraction,block=fill}
\setbeamercolor{frametitle}{bg=cMain,fg=white}
\setbeamercolor{progress bar}{fg=cAccent,bg=cMute}
\setbeamercolor{alerted text}{fg=cAccent}
\setbeamercolor{title}{fg=cMain}

% metropolis 默认尝试 Fira Sans，未安装时 fallback 到 Latin Modern Sans。
% 这里在主题加载后显式覆盖，强制使用 macOS 上一定存在的 Helvetica Neue，
% 跟中文 PingFang 风格一致；找不到则按顺序尝试 Inter / Arial。
\IfFontExistsTF{Helvetica Neue}{%
  \setsansfont{Helvetica Neue}%
}{%
  \IfFontExistsTF{Inter}{\setsansfont{Inter}}{%
    \IfFontExistsTF{Arial}{\setsansfont{Arial}}{}}%
}

% ==== TikZ ====
\usepackage{tikz}
\usetikzlibrary{positioning,arrows.meta,fit,shapes.geometric,calc,%
  decorations.pathmorphing,backgrounds}

\tikzset{
  node-box/.style={draw=cMain,thick,rounded corners=2pt,
    minimum height=8mm,minimum width=18mm,align=center,font=\small},
  node-state/.style={draw=cMain,thick,circle,minimum size=14mm,
    align=center,font=\small},
  node-ok/.style={draw=cOK,thick,fill=cOK!10},
  node-warn/.style={draw=cWarn,thick,fill=cWarn!15},
  node-bad/.style={draw=cBad,thick,fill=cBad!10},
  % 色盲友好的状态机三态：蓝 (正常/Closed) / 橙 (告警/Open) / 灰 (中间/HalfOpen)
  % 避免 red-green 二元区分，改用 hue + lightness 双通道。
  node-state-normal/.style={draw=cMain,thick,fill=cMain!12},
  node-state-alert/.style={draw=cAccent,thick,fill=cAccent!18},
  node-state-transition/.style={draw=cGray,thick,fill=cGray!18},
  arrow/.style={-{Stealth[length=2mm]},thick,draw=cMain},
  arrow-async/.style={-{Stealth[length=2mm]},thick,dashed,draw=cGray},
  arrow-bad/.style={-{Stealth[length=2mm]},thick,draw=cBad},
  arrow-alert/.style={-{Stealth[length=2mm]},thick,draw=cAccent},
  lifeline/.style={dashed,thick,draw=cGray},
}

% ==== 代码 ====
\usepackage{listings}
\lstdefinelanguage{Go}{
  morekeywords={break,case,chan,const,continue,default,defer,else,
    fallthrough,for,func,go,goto,if,import,interface,map,package,range,
    return,select,struct,switch,type,var,nil,true,false,iota},
  morekeywords=[2]{string,int,int64,uint,uint64,bool,byte,error,context},
  sensitive=true,
  morecomment=[l]//,morecomment=[s]{/*}{*/},
  morestring=[b]",morestring=[b]`,
}
\lstdefinelanguage{LuaRedis}{
  morekeywords={and,break,do,else,elseif,end,false,for,function,if,in,
    local,nil,not,or,repeat,return,then,true,until,while},
  morekeywords=[2]{redis,call,KEYS,ARGV,tonumber,HGET,HSET,GET,SET,
    DECR,DECRBY,INCR,INCRBY,EXPIRE,PEXPIRE,ZADD,ZCARD,ZREMRANGEBYSCORE,
    EVAL,EVALSHA},
  sensitive=true,
  morecomment=[l]{--},
  morestring=[b]",morestring=[b]',
}

\lstdefinestyle{base}{
  basicstyle=\ttfamily\scriptsize,
  numbers=left,numberstyle=\tiny\color{cGray},
  keywordstyle=\color{cMain}\bfseries,
  keywordstyle=[2]\color{cAccent},
  stringstyle=\color{cOK},
  commentstyle=\color{cGray}\itshape,
  backgroundcolor=\color{cMute},
  frame=single,framerule=0pt,
  showstringspaces=false,
  breaklines=true,
  tabsize=2,
  columns=fullflexible,
}
\lstset{style=base}

% ==== 三线表与附录页码 ====
\usepackage{booktabs}
\usepackage{appendixnumberbeamer}

% ==== 页脚来源标注 ====
\newcommand{\srcnote}[1]{{\color{cGray}\scriptsize\itshape #1}}

% 关键论点框
\newcommand{\keypoint}[1]{%
  \begin{tcolorbox}[colback=cMute,colframe=cMain,boxrule=0.5pt,arc=2pt]%
  #1\end{tcolorbox}}


\title{订单履约业务全景}
\subtitle{从一笔订单的生与死，到下游事件总线的回声}
\author{gomall · 业务实战}
\date{}

\begin{document}

\maketitle

\begin{frame}{今日大纲}
\tableofcontents
\end{frame}

% ============== 1. 业务背景 ==============
\section{一笔订单一生要走过的路}

\begin{frame}{为什么订单要做状态机}
\begin{itemize}
  \item 订单不是一行 DB 记录，是一段\textbf{有时间跨度的契约}：用户、商家、平台、第三方支付都在里面有义务。
  \item 任意一段时间内，订单只能处于\textbf{一个明确的业务态}：未支付、待发货、配送中、已收货、已取消、已退款。
  \item 没有状态机就没有"什么时候才能发货" / "什么时候允许退款" / "什么时候释放库存"的语义边界。
\end{itemize}
\vspace{2mm}
\keypoint{在 gomall 里，状态机 = \texttt{consts/order.go} 里那一组 iota 常量；它定义了所有合法的"现在"。}
\srcnote{consts/order.go:3--16}
\end{frame}

\begin{frame}[fragile]{gomall 的订单状态常量}
\begin{lstlisting}[language=Go,firstnumber=3,
  caption={\srcnote{consts/order.go:3--16}}]
const (
    OrderTypeUnPaid = iota + 1     // 1 未支付
    OrderTypePendingShipping       // 2 已支付 待发货
    OrderTypeShipping              // 3 已发货 待收货
    OrderTypeReceipt               // 4 已收货 完结
)
const (
    Unknown = iota
    UnPaid                         // 创建时的初始态
    Paid
    Cancelled
    Refunded
)
\end{lstlisting}
\small 两组常量分别覆盖"主流程时序"与"取消/退款旁路"。订单创建时写入的是 \texttt{consts.UnPaid}（\texttt{service/order.go:52}）。
\end{frame}

\begin{frame}{业务态四主线 + 两旁路}
\begin{tabular}{@{}llll@{}}
\toprule
状态 & 含义 & 触发者 & 时长上限 \\
\midrule
UnPaid & 未支付，已占住库存 & 用户下单 & \textbf{15--30 分钟}（关单） \\
PendingShipping & 已支付，待发货 & 用户付款 & 商家 48 小时内发货 \\
Shipping & 已发货，配送中 & 商家点发货 & 物流约 1--7 天 \\
Receipt & 已收货，交易完结 & 用户确认 / 系统超时 & 终态 \\
\midrule
Cancelled & 用户取消 / 超时关单 & 用户 / 平台调度 & 任意时点（仅未支付）\\
Refunded & 已退款 & 平台 / 商家审核 & 完结后 7--15 日内 \\
\bottomrule
\end{tabular}
\vspace{2mm}
\small 时长上限来自电商行业惯例（淘宝 / 京东公开规则），gomall 把"未支付关单"硬编码为 15 分钟（见 \texttt{service/order.go:89} 的 zset score）和 30 分钟（见 \texttt{repository/rabbitmq/order\_delay.go:22}）两套互补机制。
\srcnote{service/order.go:88--98 + repository/rabbitmq/order\_delay.go:22}
\end{frame}

\begin{frame}{订单状态机全景图}
\begin{tikzpicture}[node distance=10mm and 16mm,font=\scriptsize]
  \node[node-state,node-state-normal] (up) {UnPaid};
  \node[node-state,node-state-normal,right=of up] (ps) {Pending\\Shipping};
  \node[node-state,node-state-normal,right=of ps] (sh) {Shipping};
  \node[node-state,node-state-normal,right=of sh] (rc) {Receipt};
  \node[node-state,node-state-alert,below=of up] (cc) {Cancelled};
  \node[node-state,node-state-alert,below=of rc] (rf) {Refunded};

  \draw[arrow] (up) -- node[above,font=\tiny]{PayDown} (ps);
  \draw[arrow] (ps) -- node[above,font=\tiny]{商家发货} (sh);
  \draw[arrow] (sh) -- node[above,font=\tiny]{用户确认} (rc);

  \draw[arrow-alert] (up) -- node[left,font=\tiny]{超时 / 主动取消} (cc);
  \draw[arrow-alert] (rc) -- node[right,font=\tiny]{售后退款} (rf);
  \draw[arrow-alert,bend right=15] (ps) to node[below right,font=\tiny]{未发货退款} (rf);
\end{tikzpicture}
\vspace{2mm}
\small 主流程是一条单向链；Cancelled 只能从 UnPaid 进入；Refunded 可从 PendingShipping / Receipt 进入。
\srcnote{service/payment.go:146 (UnPaid$\to$PendingShipping) + service/order\_cancel.go:32 (UnPaid$\to$Cancelled)}
\end{frame}

% ============== 2. 库存预扣 ==============
\section{为什么下单要预扣库存而不是直接扣 DB}

\begin{frame}{业务问题：库存到底什么时候"少一件"}
\begin{itemize}
  \item 直觉答案：用户下单时少一件。
  \item 但下单 $\ne$ 付款。用户挂着购物车不付款的比例非常高（电商业内通用统计在 60--70\%）。
  \item 如果下单立即扣 DB \texttt{product.num}：
  \begin{itemize}
    \item 大量未支付订单 $\to$ DB 库存被假占用 $\to$ 真正想买的用户看到"无货"
    \item 取消时必须把库存\textbf{加回去}，一旦"取消路径"和"下单路径"不对称就出 bug
  \end{itemize}
\end{itemize}
\srcnote{stressTest/REPORT.md \S 已修：库存虚高 + 历史 PR \#51 / \#58}
\end{frame}

\begin{frame}{gomall 的设计：两个桶}
\begin{tikzpicture}[node distance=28mm]
  \node[node-box,minimum width=34mm,minimum height=18mm,node-ok]
    (av) {available\\\texttt{stock:available:\{pid\}}};
  \node[node-box,minimum width=34mm,minimum height=18mm,node-warn,right=of av]
    (rv) {reserved\\\texttt{stock:reserved:\{pid\}}};
  \node[node-box,minimum width=34mm,minimum height=18mm,below=18mm of rv,node-bad]
    (sold) {sold\\(DB \texttt{product.num} 减少)};
  \draw[arrow,bend left=18] (av) to node[above,font=\tiny]{reserve(n) 下单} (rv);
  \draw[arrow,bend left=18] (rv) to node[below,font=\tiny]{release(n) 取消/超时} (av);
  \draw[arrow] (rv) -- node[right,font=\tiny]{commit(n) 支付成功} (sold);
\end{tikzpicture}
\vspace{1mm}
\keypoint{不变量：\texttt{available + reserved + sold = 初始库存}。任何一步只在两个桶间挪，不会凭空多/少。}
\srcnote{repository/cache/inventory.go:11--14 + 24--75}
\end{frame}

\begin{frame}[fragile]{OrderCreate 里这套桶怎么用}
\begin{lstlisting}[language=Go,firstnumber=57,
  caption={\srcnote{service/order.go:57--86}}]
// 1) Redis 预扣库存: available -> reserved，失败直接拒单
if err = cache.ReserveStock(ctx, req.ProductID, int64(req.Num)); err != nil {
    return nil, err
}
// 2) 同事务写订单 + 写 outbox 事件
err = dao.NewDBClient(ctx).Transaction(func(tx *gorm.DB) error {
    if e := dao.NewOrderDaoByDB(tx).CreateOrder(order); e != nil {
        return e
    }
    return dao.NewOutboxDaoByDB(tx).Insert(...)
})
if err != nil {
    // 3) 事务失败：释放刚扣的预占库存，回到 available
    cache.ReleaseReservation(ctx, req.ProductID, int64(req.Num))
}
\end{lstlisting}
\small 三步顺序：预扣 $\to$ 写订单 $\to$ 失败回退。预扣失败直接拒单；订单写失败回退库存 —— 状态机闭环。
\end{frame}

\begin{frame}{压测验证：500 并发抢 100 件零超卖}
\begin{tabular}{@{}lr@{}}
\toprule
项 & 数值 \\
\midrule
初始 available 桶 & 100 \\
并发 goroutine 数 & 500 \\
每个 goroutine 试图 reserve & 1 件 \\
\midrule
成功 reserve 数 (\texttt{success}) & \textbf{100} \\
available 桶最终值 & \textbf{0} \\
reserved 桶最终值 & \textbf{100} \\
超卖次数 & \textbf{0} \\
\bottomrule
\end{tabular}
\vspace{2mm}
\small 单条 Lua 脚本里 \texttt{GET}+\texttt{DECRBY}+\texttt{INCRBY} 原子执行，500 个 goroutine 排队跑，正好抢光 100 件。
\srcnote{repository/cache/inventory\_test.go:52--82 \texttt{TestInventory\_NoOversellUnderConcurrency}}
\end{frame}

% ============== 3. 事件流 ==============
\section{一笔订单的事件流}

\begin{frame}{时间轴：从下单点击到完结}
\begin{tikzpicture}[node distance=4mm and 6mm,font=\tiny,
    every node/.append style={minimum width=15mm,minimum height=7mm}]
  \node[node-box,node-ok] (a) {点下单};
  \node[node-box,node-warn,right=of a] (b) {Reserve\\预扣};
  \node[node-box,right=of b] (c) {INSERT\\order};
  \node[node-box,right=of c] (d) {Outbox\\order.created};
  \node[node-box,below=4mm of d] (e) {用户付款};
  \node[node-box,left=of e] (f) {Commit\\reserved};
  \node[node-box,left=of f] (g) {Outbox\\order.paid};
  \node[node-box,left=of g] (h) {商家发货};
  \node[node-box,below=4mm of h] (i) {用户收货};
  \node[node-box,node-ok,right=of i] (j) {完结};
  \draw[arrow] (a) -- (b);
  \draw[arrow] (b) -- (c);
  \draw[arrow] (c) -- (d);
  \draw[arrow,dashed] (d) -- (e);
  \draw[arrow] (e) -- (f);
  \draw[arrow] (f) -- (g);
  \draw[arrow,dashed] (g) -- (h);
  \draw[arrow,dashed] (h) -- (i);
  \draw[arrow] (i) -- (j);
\end{tikzpicture}
\vspace{2mm}
\small 实线 = 同步发生在请求里；虚线 = 异步等待（用户操作 / 物流时间）。Outbox 把异步段的开始时刻\textbf{记录下来}。
\end{frame}

\begin{frame}{order.paid 事件：动作分裂为多条业务支线}
\begin{itemize}
  \item \textbf{库存正式扣}：\texttt{commit(reserved)} 在支付事务里同步做了；事件不再负责扣库存，但负责\textbf{通知}。
  \item \textbf{商家通知}：发短信 / IM，"您有新订单"。
  \item \textbf{推荐系统}：用户买了什么 $\to$ 用户画像更新 $\to$ 召回池调整。
  \item \textbf{对账系统}：把 order\_id + amount 入对账明细，T+1 跟支付通道核对。
  \item \textbf{BI}：实时 GMV / 客单价 / 转化率。
  \item \textbf{营销返券}：满足规则的订单触发返券、积分、等级。
\end{itemize}
\srcnote{service/payment.go:174--184 (写 outbox order.paid) + service/events/events.go:17--23}
\end{frame}

\begin{frame}{order.cancelled 事件：补偿与通知}
\begin{itemize}
  \item \textbf{库存释放}：\texttt{ReleaseReservation()} 在写完 outbox 后同步执行（\texttt{service/order\_cancel.go:59}）—— 桶状态机回退。
  \item \textbf{商家通知}：发"取消"提醒，让商家更新待发货列表。
  \item \textbf{营销回滚}：如果下单时发过预占券 / 满减锁定，要释放回去。
  \item \textbf{BI}：转化漏斗里这一单标记为"流失"，纳入流失分析。
  \item \textbf{风控}：高频取消 + 短时间内多账号 $\to$ 套券 / 刷单嫌疑。
\end{itemize}
\vspace{2mm}
\srcnote{service/order\_cancel.go:30--62 + service/events/events.go:25--32}
\end{frame}

% ============== 4. Outbox 业务侧 ==============
\section{Outbox 在业务侧的角色}

\begin{frame}{Outbox 把"一次业务"变成"一束事件"}
\begin{tikzpicture}[node distance=6mm and 14mm,font=\scriptsize]
  \node[node-box,minimum width=24mm,minimum height=12mm,node-warn] (tx)
    {同一个 DB 事务\\order + outbox};
  \node[node-box,right=14mm of tx] (pub) {publisher};
  \node[node-box,right=of pub,minimum width=20mm] (mq) {RabbitMQ\\\texttt{domain.events}};
  \node[node-box,right=14mm of mq.north east,anchor=west,node-ok] (sub1) {对账};
  \node[node-box,below=of sub1,node-ok] (sub2) {BI};
  \node[node-box,below=of sub2,node-ok] (sub3) {风控};
  \node[node-box,below=of sub3,node-ok] (sub4) {推荐};
  \draw[arrow] (tx) -- (pub);
  \draw[arrow] (pub) -- (mq);
  \draw[arrow-async] (mq) -- (sub1);
  \draw[arrow-async] (mq) -- (sub2);
  \draw[arrow-async] (mq) -- (sub3);
  \draw[arrow-async] (mq) -- (sub4);
\end{tikzpicture}
\vspace{2mm}
\small 业务侧不需要"逐一通知"。它只写一行 outbox，下游愿意来订就来订 —— \textbf{业务和增长解耦}。
\srcnote{repository/db/model/outbox.go:17--37}
\end{frame}

\begin{frame}[fragile]{outbox\_event 表结构再看一眼}
\begin{lstlisting}[language=Go,firstnumber=24,
  caption={\srcnote{repository/db/model/outbox.go:24--35}}]
type OutboxEvent struct {
    gorm.Model
    AggregateType string  // "order" / "product"
    AggregateID   uint    // 业务实体 id
    EventType     string  // "OrderCreated" / "OrderCancelled"
    RoutingKey    string  // "order.created"
    Payload       string  // JSON body (events.OrderCreated 等)
    Status        int     // 1 pending / 2 sent / 3 dead
    Attempts      int
    NextRetryAt   time.Time
    LastError     string
}
\end{lstlisting}
\small \texttt{AggregateType}+\texttt{AggregateID} 是业务追踪根；\texttt{RoutingKey} 决定谁能订阅；\texttt{Status} 是 publisher 的视角。
\end{frame}

\begin{frame}{业务侧下游订阅者清单}
\begin{tabular}{@{}lll@{}}
\toprule
下游 & 关注事件 & 业务价值 \\
\midrule
对账 & order.paid / order.refunded & 资金流闭环，T+1 跟支付通道核对 \\
BI & 全量 order.* & GMV / 客单 / 转化漏斗 \\
风控 & order.created / cancelled & 异常订单识别 \\
推荐 & order.paid & 行为打点 $\to$ 召回排序 \\
营销 & order.created / paid & 返券 / 满减判断 \\
商家工作台 & order.paid / shipped & 待办列表 \\
ES 索引 & product.changed / order.paid & 搜索增量索引 \\
\bottomrule
\end{tabular}
\vspace{2mm}
\small 真实电商每条 routing key 背后至少 3--5 个订阅者；gomall 现网先接了 ES 索引（\texttt{service/search/indexer.go}）。
\srcnote{service/search/indexer.go:28 (Consume) + repository/rabbitmq/domain.go:10}
\end{frame}

% ============== 5. 事故复盘 ==============
\section{真实事故：库存虚高 bug 业务复盘}

\begin{frame}[fragile]{现象：库存数字一路膨胀}
\begin{itemize}
  \item 商家投诉："我商品明明只有 100 件，怎么后台显示 130 / 160 / 200？"
  \item 售前页面显示有货，用户下单付钱，售后却"无货发不出"。
  \item 客服侧不断收到"已付款但商家说没货"的工单。
\end{itemize}
\vspace{2mm}
\small 这是一个典型的\textbf{业务表象 $\to$ 代码根因}的链条：现象在客服系统里，根因在两个 Go 函数里。
\srcnote{stressTest/REPORT.md \S 已修：库存虚高 + service/order\_task.go:33--35 \texttt{RollbackStock}}
\end{frame}

\begin{frame}[fragile]{根因：一对不对称的孪生函数}
\begin{lstlisting}[language=Go,basicstyle=\ttfamily\scriptsize]
// 旧版 OrderCreate (PR #58 前)
// -> 没有任何 product.Num -= n 的语句
//    库存不动，只写订单

// 旧版超时关单 (service/order_task.go:33)
//    txProductDao.RollbackStock(order.ProductID, order.Num)
//    -> UPDATE product SET num = num + n
\end{lstlisting}
\vspace{2mm}
\begin{itemize}
  \item 下单不扣，关单却"还回去"。
  \item 每个超时未支付订单 = 库存凭空 +N。
  \item 几个月累计下来，热门商品的 \texttt{product.num} 离真实值越来越远。
\end{itemize}
\srcnote{service/order\_task.go:33 + 历史 PR \#51 引入延迟关单时埋下伏笔}
\end{frame}

\begin{frame}{业务推演：单商品 100 件，跑一个月}
\begin{tikzpicture}[node distance=8mm and 14mm,font=\scriptsize]
  \node[node-box,node-ok] (init) {Day 0\\product.num = 100};
  \node[node-box,below right=of init,node-warn] (d10) {Day 10\\超时未付 30 单\\$\to$ +30};
  \node[node-box,right=of d10,node-bad] (d20) {Day 20\\再来 40 单\\$\to$ 170};
  \node[node-box,below=of d20,node-bad] (d30) {Day 30\\已售 100 件\\但显示 220};
  \draw[arrow] (init) -- (d10);
  \draw[arrow] (d10) -- (d20);
  \draw[arrow] (d20) -- (d30);
\end{tikzpicture}
\vspace{2mm}
\small "下单不扣 + 取消加回"这对不对称行为，配合电商常态化的 60\%+ 弃单率，几乎注定会爆。
\end{frame}

\begin{frame}{修法：把库存状态机显式化}
\begin{tabular}{@{}lll@{}}
\toprule
方案 & 改动量 & 评价 \\
\midrule
A. 加一列 \texttt{reserved\_num} & 小 & 治标，状态机仍隐式 \\
B. 双桶 + reserve/commit/release & 中 & \textbf{推荐}：显式三态、闭环对称 \\
\bottomrule
\end{tabular}
\vspace{2mm}
\begin{itemize}
  \item 选 B（即 Deck 04 的双桶设计）。
  \item 下单 = reserve（available $\to$ reserved），不动 DB。
  \item 付款 = commit（reserved $\to$ 真扣 DB）。
  \item 取消 = release（reserved $\to$ available），\textbf{不再回写 DB}。
\end{itemize}
\srcnote{service/order\_cancel.go:14--18 关键注释：不再回写 product.Num}
\end{frame}

% ============== 6. 性能与业界 ==============
\section{性能与业界对照}

\begin{frame}{订单列表压测：游标分页的代差碾压}
\begin{tabular}{@{}lrr@{}}
\toprule
指标 & 优化版（游标 + 缓存） & 旧版（OFFSET 深分页） \\
\midrule
吞吐 (RPS) & \textbf{58,406} & 8.3 \\
p50 延迟 & 1.04 ms & 10.5 s \\
p95 延迟 & \textbf{5 ms} & 15.95 s \\
最大延迟 & 348 ms & 16.29 s \\
错误率 & 0\% & 0\% \\
\bottomrule
\end{tabular}
\vspace{2mm}
\small 同一个业务"我的订单列表"，深分页时 p95 接近 16 秒 —— 这在真实业务里等于直接超时；用游标 + Redis 缓存的优化版 p95 拉到 5 ms，相差\textbf{约 3,200 倍}。
\srcnote{stressTest/REPORT.md \S 2. 订单列表：游标分页的代差碾压}
\end{frame}

\begin{frame}{业界 framing：电商订单系统怎么设计}
\begin{itemize}
  \item \textbf{阿里}：交易中心把订单建模成"状态机 + 事件流"，公开技术博客多次提到"扁平化状态机"+"领域事件总线"。
  \item \textbf{京东}：履约链路按"接单 / 履约 / 完结"三大域，每个域独立状态机，通过事件解耦。QCon 公开分享强调"补偿事务 + 异步消息"是底线设计。
  \item \textbf{Saga 模式}：业界电商对长事务的事实标准 —— gomall 用 Outbox + 协同式 Saga 是同一族思路。
\end{itemize}
\vspace{2mm}
\begin{tabular}{@{}lll@{}}
\toprule
维度 & gomall & 大型电商 \\
\midrule
状态数 & 4 主 + 2 旁 & 数十个细分 \\
事件总线 & 单 RabbitMQ exchange & 多级 / 跨域 / 多机房 \\
补偿事务 & 协同式 Saga + 释放库存 & 编排式 Saga + 人工兜底 \\
\bottomrule
\end{tabular}
\vspace{1mm}\\
\small \textbf{原理一致，工程量天差地别}。这里只列方向不引数字，避免不可考的"百万 TPS"之类。
\end{frame}

% ============== Q&A ==============
\appendix

\begin{frame}{面试 Q\&A（上）}
\textbf{Q1}：为什么下单不直接扣 DB？\\
\hspace{4mm}\small A：DB 库存 $\ne$ 售前可见库存。下单只是占用，付款才是消耗。两者混淆会导致"超时单 + 反向补偿"的不对称 bug。预扣（reserved 桶）就是把"占用"这个时间窗口显式建模。

\vspace{2mm}
\textbf{Q2}：reserved 桶里的库存，售前要怎么显示？\\
\hspace{4mm}\small A：售前页面读 \texttt{available} 桶 —— reserved 对售前不可见。\\\texttt{available + reserved + sold = 初始水位} 这个不变量保证你怎么算都不会超。

\vspace{2mm}
\textbf{Q3}：为什么不直接 DB 加锁 \texttt{SELECT FOR UPDATE}？\\
\hspace{4mm}\small A：可以，但每次下单都锁行 + 等锁排队，吞吐会塌。Redis Lua 是\textbf{单线程原子}，性能差一个数量级。本项目压测 500 并发零超卖（\texttt{repository/cache/inventory\_test.go:52--82}）。
\end{frame}

\begin{frame}{面试 Q\&A（下）}
\textbf{Q4}：order.paid 事件的下游有哪些？\\
\hspace{4mm}\small A：对账、商家通知、推荐、BI、营销返券、ES 索引。Outbox + RabbitMQ 把它们都解耦 —— 业务侧只写一行 \texttt{outbox.Insert}。

\vspace{2mm}
\textbf{Q5}：库存虚高事故你怎么排查？\\
\hspace{4mm}\small A：从业务现象（商家投诉、售前售后不一致）反推：库存只增不减 $\to$ 找所有 \texttt{product.Num +=} 的地方 $\to$ 发现 \texttt{RollbackStock}，再找有没有对称的 \texttt{-=} $\to$ 发现 OrderCreate 根本没扣 $\to$ 根因是状态机隐式 / 路径不对称。修法：显式化双桶状态机。

\vspace{2mm}
\textbf{反追问}：如果未付款订单有几十万，关单怎么扛？\\
\hspace{4mm}\small A：本项目用 RabbitMQ 死信队列做延迟关单（\texttt{repository/rabbitmq/order\_delay.go:22} 默认 30 分钟）+ 定时任务兜底（\texttt{service/order\_task.go} 扫超过 15 分钟的）。两条路径都做了幂等（\texttt{CloseOrderWithCheck}），重复触发 no-op。
\end{frame}

\begin{frame}{代码位置一览}
\begin{description}
  \item[订单状态常量]\srcnote{consts/order.go:3--16}
  \item[OrderCreate + 预扣 + Outbox]\srcnote{service/order.go:40--102}
  \item[PayDown + 状态转移 + Commit]\srcnote{service/payment.go:35--200}
  \item[CancelUnpaidOrder + Release]\srcnote{service/order\_cancel.go:14--63}
  \item[双桶库存状态机]\srcnote{repository/cache/inventory.go:16--138}
  \item[Outbox 表 + 事件类型]\srcnote{repository/db/model/outbox.go:17--37}
  \item[事件 payload]\srcnote{service/events/events.go:9--38}
  \item[延迟关单拓扑]\srcnote{repository/rabbitmq/order\_delay.go:13--105}
  \item[500 并发库存测试]\srcnote{repository/cache/inventory\_test.go:52--82}
\end{description}
\vspace{1mm}
\keypoint{业务复盘的关键引用：\texttt{service/order\_cancel.go:14--18} 注释里写明"不再回写 product.Num"。}
\end{frame}

\end{document}
